% Оборот титульного листа: аннотация, адресат, как устроены примеры.
\pagestyle{empty}
\sffamily
\setlength{\parindent}{0pt}
\vspace*{\fill}

{\large\bfseries\color{GeoskanDarkBlue} Корягин С.\,А.}\par
\vspace{4pt}
{\color{GeoskanDarkBlue} Программирование дронов и создание автономных беспилотных миссий: учебное пособие по курсу «Программные средства БПЛА». --- Кудрово, \the\year. --- \pageref{LastPage}~с.}\par
\vspace{1cm}

{\bfseries\color{GeoskanDarkBlue} Аннотация}\par
\vspace{4pt}
Пособие вводит в программирование образовательного квадрокоптера «Геоскан Пионер» на языках Lua и Python. Первая часть --- язык Lua, светодиодная индикация, таймеры и событийная модель автопилота, автономные миссии по точкам и параметрическим траекториям. Вторая часть --- Python и библиотека \texttt{pioneer\_sdk}: подключение, полёт, камера, визуальные маркеры ArUco, трекинг и групповой полёт. Третья часть --- профессиональная навигация: калибровка камеры, карты маркеров, фильтрация измерений. Все программы запускаются в веб-симуляторе \texttt{simulator.sakoryagin.ru} с тем же API, что у настоящего дрона.\par
\vspace{0.8cm}

{\bfseries\color{GeoskanDarkBlue} Для кого}\par
\vspace{4pt}
Для школьников и студентов, начинающих программировать роботов, и педагогов дополнительного образования. Предварительной подготовки по программированию не требуется; математика --- в объёме школьной тригонометрии.\par
\vspace{0.8cm}

{\bfseries\color{GeoskanDarkBlue} Примеры и проверка}\par
\vspace{4pt}
Исходный код всех примеров --- в папке \texttt{code/} рядом с книгой; каждый файл подписан в тексте. Примеры автоматически прогоняются в симуляторе: взлетают, проходят точки и садятся. Текст примеров можно копировать прямо из PDF.\par
\vspace{0.8cm}

{\footnotesize\color{GeoskanMuted} Версия \theversion\ · сборка \buildtimestamp\par}
\vspace*{\fill}
\clearpage
\pagestyle{scrheadings}
